\chapter{Literature Review}
\label{ch:literature}

The detection of malicious, scam, and phishing websites has been an active area of research for over two decades, evolving alongside the sophistication of online threats. Researchers have explored a wide range of signals — URL structure, domain metadata, page content, network behavior, and threat intelligence feeds — to build reliable detection systems. The following studies highlight the core ideas, methodologies, and technologies that informed and shaped the development of ProblySUS.

\section{URL-Based Phishing Detection Using Lexical Features}

Early research in phishing detection established that a URL's lexical structure alone carries significant predictive power. Studies demonstrated that features such as the use of raw IP addresses as hostnames, excessive hyphens in domain names, the presence of brand names embedded in subdomains, the use of free or abused top-level domains (e.g. \texttt{.tk}, \texttt{.ml}, \texttt{.xyz}), and phishing-related keywords in the URL path (e.g. \texttt{login}, \texttt{verify}, \texttt{paypal}, \texttt{secure}) are strong and consistent indicators of fraudulent intent. Machine learning classifiers trained on these features achieved high detection accuracy while remaining computationally lightweight and requiring no page access.

This directly informs ProblySUS's \texttt{pattern\_checker} module, which performs zero-network, instant lexical inspection of every submitted URL before any heavier analysis begins. By flagging IP-based hostnames, abused TLDs, hyphen counts, and keyword patterns, the module acts as a fast pre-filter that can assign risk points within milliseconds, forming the foundation of the composite scoring pipeline.

\section{Domain Registration and WHOIS Intelligence in Fraud Detection}

Research has consistently identified domain age as one of the strongest individual signals for phishing and scam detection. Fraudulent websites are typically registered shortly before being used in an attack campaign, with the majority of phishing domains being abandoned within 30 to 90 days of registration. Studies leveraging WHOIS registration data showed that newly registered domains (under 30 days old) carry a disproportionately high likelihood of being associated with malicious activity, and that registrar patterns, registration country, and WHOIS privacy shielding also correlate with fraud risk.

ProblySUS's \texttt{whois\_checker} module directly applies this insight by querying the WHOIS registry for each submitted domain's creation date and computing its age in days. Domains under one week old are assigned the maximum domain-age penalty, those under one month receive a high penalty, and those under three months receive a moderate penalty — thresholds calibrated from the literature to minimize false positives against legitimately new but benign sites.

\section{Threat Intelligence Feeds and Blacklist-Based Detection}

Community-maintained threat intelligence feeds represent one of the most reliable and widely adopted mechanisms for malicious URL detection. Research evaluating platforms such as URLhaus, PhishTank, and Google Safe Browsing demonstrated that blacklist-based detection achieves near-perfect precision for known threats, with the primary limitation being coverage of newly emerging URLs that have not yet been reported. Studies recommended combining blacklist lookups with complementary heuristic signals to bridge this gap, producing a hybrid detection approach that handles both known and zero-day threats.

ProblySUS integrates URLhaus — a community-maintained, freely available threat feed — as a hard-override detection layer. The \texttt{blacklist\_checker} and \texttt{updater} modules download the URLhaus ZIP feed, extract the CSV, and maintain a local indexed copy that is auto-refreshed every 24 hours. A confirmed blacklist hit returns an instant score of 100 (Fraudulent), bypassing all other analysis, while the remaining modules handle threats not yet listed in the feed.

\section{Content and Behavioral Signals in Phishing Page Detection}

Studies on content-based phishing detection highlighted that fraudulent pages consistently exhibit a distinct set of behavioral and structural patterns: urgency and pressure language (``Act now'', ``Your account will be suspended''), absence of legitimate trust pages (Privacy Policy, Contact, Terms of Service), requests for sensitive credentials or payment information over unencrypted connections, and unusually high external link ratios. Complementary research on JavaScript-driven attacks demonstrated that many modern phishing pages use client-side redirects, dynamic resource loading, and obfuscated scripts to evade static analysis — underlining the need for runtime browser-based inspection.

These findings motivated two distinct modules in ProblySUS. The \texttt{content\_checker} module fetches live page HTML and scans for urgency phrases, sensitive data forms, trust indicators, and reading complexity. The \texttt{behavior\_checker} module addresses the JavaScript evasion challenge by launching a headless Chromium browser via Playwright, which executes JavaScript exactly as a real browser would — intercepting all HTTP responses to detect redirect chains and capturing all outgoing network requests to identify suspicious runtime domains that no static scan could observe.

\section{Network Footprint Analysis and Privacy Profiling of Websites}

Recent research extended phishing and fraud detection beyond the primary domain to examine a site's broader network footprint — the external services, CDN providers, advertising networks, and analytics platforms it contacts. Studies found that fraudulent sites disproportionately contact domains with suspicious TLDs, IP-based hosts, and excessive numbers of external dependencies, while legitimate sites typically rely on a small set of well-known infrastructure providers (Cloudflare, AWS, Google, Fastly). Parallel work on web privacy identified browser fingerprinting — the use of Canvas, WebGL, AudioContext, and WebRTC APIs to uniquely identify visitors — as a growing surveillance technique employed on a wide range of sites, including fraudulent ones.

ProblySUS operationalizes these findings through two dedicated modules. The \texttt{network\_checker} module classifies every domain contacted during the Playwright session into \texttt{safe\_infra}, \texttt{suspicious}, or \texttt{unknown} buckets and computes an overall network risk level (Low/Medium/High). The \texttt{privacy\_checker} module independently measures cookie tracking intensity, third-party script count, and the presence of fingerprinting API calls, assigning a privacy grade from Good to Invasive — providing users with a holistic view of a site's data collection practices beyond its direct fraud risk.

Together, these research directions form the theoretical and technical foundation of ProblySUS, collectively motivating its multi-signal, layered approach to scam detection. The combination of lexical URL analysis, WHOIS domain intelligence, community threat feeds, live content and behavior inspection, and network privacy profiling produces a detection system that is both comprehensive in coverage and explainable in output — addressing the core limitations of any single-signal approach and providing everyday users with a transparent, actionable verdict on the trustworthiness of any website.